% archon:future-covers MR4689373HigherSiegelWeilFormulaForUnitaryGroupsTheNonSingularTerms/Assembly.lean
\chapter{Assembly: Proof of the Main Theorem}
\label{chap:assembly}

This chapter formalizes Phase~P3d of the project: combining the results of Chapters
\ref{chap:perverse-sheaves-K} and~\ref{chap:wd-representations} to prove Theorem~1.1
(Theorem~\ref{thm:higher-siegel-weil} of Chapter~\ref{chap:overview}). The proof is a
chain of equalities — differentiation lemma, perverse sheaf isomorphism, and Frobenius
trace on the intersection side — each previously established. We follow \S\S15--16 of
Feng--Yun--Zhang [MR4689373].

\section{Matching of sheaves}

\subsection{Proof of the main theorem}

\begin{lemma}[Differentiation lemma]
  \label{lem:differentiate-q}
  \lean{HSW.differentiateQ}
  \uses{thm:trace-Eis, def:normalized-fourier-coeff, def:central-value-d, def:K-Eis}
  % SOURCE: FYZ-HigherSiegelWeil.tex, §16.1–16.2
  % SOURCE QUOTE: "\(\frac{1}{(\log q)^r}(\dds)^r|_{s=0}(q^{ds}\wt E_a(m(\cE),s,\Phi))
  %   =\Tr(\Frob_q^{(r)},(\cK_d^\Eis)_{\bar\cQ})\)"
  \textit{Source: Feng--Yun--Zhang [MR4689373], \S16.1--16.2.}
  Let \(d \geq 0\) and \(\cQ \in \Herm_{2d}(\mathbf{F}_q)\). Using the graded
  Frobenius trace formula~\ref{thm:trace-Eis}, the \(r\)-th normalized logarithmic
  \(s\)-derivative at \(s = 0\) satisfies:
  \begin{align*}
    &\frac{1}{(\log q)^r}
    \Bigl(\frac{d}{ds}\Bigr)^r\!\Bigg|_{s=0}
    \!\bigl(q^{ds}\,\widetilde{E}_a(m(\cE), s, \Phi)\bigr)
    \\
    &\quad=
    \frac{1}{(\log q)^r}
    \Bigl(\frac{d}{ds}\Bigr)^r\!\Bigg|_{s=0}
    \Biggl(
      q^{ds} \sum_{i=0}^d
      \Tr\!\bigl(\Frob_q,\; (\cK^{\mathrm{Eis}}_{d,i})_{\overline{\cQ}}\bigr)\,
      q^{-2is}
    \Biggr)
    \\
    &\quad=
    \sum_{i=0}^d
    \Tr\!\bigl(\Frob_q,\; (\cK^{\mathrm{Eis}}_{d,i})_{\overline{\cQ}}\bigr)\cdot
    (d - 2i)^r.
  \end{align*}
  The last step uses the identity
  \(\frac{1}{(\log q)^r}(\frac{d}{ds})^r|_{s=0}(q^{(d-2i)s}) = (d-2i)^r\).
\end{lemma}

\begin{proof}
  \uses{thm:trace-Eis, def:normalized-fourier-coeff}
  Substitute the graded trace formula (Theorem~\ref{thm:trace-Eis}) into
  \(q^{ds}\widetilde{E}_a(m(\cE),s,\Phi)\):
  \[
    q^{ds}\widetilde{E}_a(m(\cE),s,\Phi)
    = \sum_{i=0}^d
    \Tr\!\bigl(\Frob_q,\; (\cK^{\mathrm{Eis}}_{d,i})_{\overline{\cQ}}\bigr)\,
    q^{(d-2i)s}.
  \]
  Now differentiate term by term: \((\frac{d}{ds})^r q^{(d-2i)s} =
  ((d-2i)\log q)^r\, q^{(d-2i)s}\), evaluated at \(s=0\) gives
  \(((d-2i)\log q)^r\). Dividing by \((\log q)^r\) yields \((d-2i)^r\) per term,
  completing the computation.
\end{proof}

\begin{theorem}[Main assembly]
  \label{thm:main-assembly}
  \lean{HSW.mainAssembly}
  \uses{lem:differentiate-q, thm:trace-Int, cor:K-Eis-iso-K-Int,
        thm:higher-siegel-weil, def:degree-special-cycle, def:virtual-class-full}
  % SOURCE: FYZ-HigherSiegelWeil.tex, §16.3
  % SOURCE QUOTE: "Combining steps (1)--(3), Theorem~\ref{th:intro} follows"
  \textit{Source: Feng--Yun--Zhang [MR4689373], \S16.3.}
  The higher Siegel--Weil formula (Theorem~\ref{thm:higher-siegel-weil}) holds:
  \[
    \frac{1}{(\log q)^r}
    \left(\frac{d}{ds}\right)^r\!\Bigg|_{s=0}
    \!\bigl(q^{d s}\,\widetilde{E}_a(m(\cE), s, \Phi)\bigr)
    \;=\; \deg[\cZ_{\cE}^r(a)].
  \]
\end{theorem}

\begin{proof}
  \uses{lem:differentiate-q, thm:trace-Int, cor:K-Eis-iso-K-Int}
  We chain three equalities:
  \begin{align*}
    &\frac{1}{(\log q)^r}
    \Bigl(\frac{d}{ds}\Bigr)^r\!\Big|_{s=0}
    \bigl(q^{ds}\widetilde{E}_a(m(\cE),s,\Phi)\bigr) \\
    &\overset{\text{Lem.~\ref{lem:differentiate-q}}}{=\!=\!=}
    \sum_{i=0}^d
    \Tr\!\bigl(\Frob_q,\; (\cK^{\mathrm{Eis}}_{d,i})_{\overline{\cQ}}\bigr)\cdot
    (d-2i)^r \\
    &\overset{\text{Cor.~\ref{cor:K-Eis-iso-K-Int}}}{=\!=\!=}
    \sum_{i=0}^d
    \Tr\!\bigl(\Frob_q,\; (\cK^{\mathrm{Int}}_{d,i})_{\overline{\cQ}}\bigr)\cdot
    (d-2i)^r \\
    &\overset{\text{Thm.~\ref{thm:trace-Int}}}{=\!=\!=}
    \deg[\cZ_{\cE}^r(a)].
  \end{align*}
  \textbf{Step~1} is Lemma~\ref{lem:differentiate-q}: differentiation converts the
  graded trace of \(\cK_d^{\mathrm{Eis}}\) at \(s = 0\) into the sum
  \(\sum_i \Tr(\cK^{\mathrm{Eis}}_{d,i}) \cdot (d-2i)^r\).

  \textbf{Step~2} uses the isomorphism \(\cK_d^{\mathrm{Eis}} \cong \cK_d^{\mathrm{Int}}\)
  (Corollary~\ref{cor:K-Eis-iso-K-Int}, Chapter~\ref{chap:wd-representations}), which
  replaces each graded stalk trace of \(\cK^{\mathrm{Eis}}_{d,i}\) with the corresponding
  trace of \(\cK^{\mathrm{Int}}_{d,i}\).

  \textbf{Step~3} applies the Frobenius trace formula for \(\cK_d^{\mathrm{Int}}\)
  (Theorem~\ref{thm:trace-Int}, Chapter~\ref{chap:perverse-sheaves-K}), which identifies
  \(\sum_i \Tr(\cK^{\mathrm{Int}}_{d,i}) \cdot (d-2i)^r\) with \(\deg[\cZ_{\cE}^r(a)]\).

  Each step is a previously proved result; the assembly is purely formal. This completes
  the proof of the main theorem.
\end{proof}
